\documentclass[11pt,a4paper]{proofarticle}
\usepackage{graphicx,xcolor,relinput}
\relinput{.}{authors_cmd}
\DeclareUnicodeCharacter{2032}{\prime}
\nocite{*}
\title{Cadre Analytique Modulaire pour l’Hypothèse de Riemann (T1$′$–T4)\\[2pt]
\large Preuve-seule, prêt-journal, compatible arXiv}
\author{\InterIAauthors}
\date{\today}
\begin{document}

\relinput{.}{docs/explication-grand-format/latex/page_de_garde_FR}

\maketitle

\section*{Préambule}
Ce document présente le cadre analytique modulaire (T1$′$–T4) reliant
la formule explicite de Riemann–Weil, les encadrements spectraux,
et la lecture interdisciplinaire proposée par le collectif InterIA.

\relinput{.}{appendix/table_des_symboles_FR}
\addcontentsline{toc}{section}{Sommaire des notations}

% --- section Itinéraires de lecture
\relinput{.}{docs/explication-grand-format/latex/00_persona_routes_FR}

% --- section Dictionnaire croisé
\relinput{.}{docs/explication-grand-format/latex/01_translation_dictionary_FR}

% --- section Storyboard T1$′$–T4
\relinput{.}{docs/explication-grand-format/latex/02_storyboard_t1t4_FR}

% --- section Micro-exemple (T2)
\relinput{.}{docs/explication-grand-format/latex/03_micro_example_t2_FR}

% --- section T2 — Weighted frame bound (full)
\relinput{.}{theorems/T2_borne_de_frame_complete_FR} % T2_weighted_frame_full

% --- section T3 — Wave–packet localization (full)
\relinput{.}{theorems/T3_localisation_paquet_ondes_FR} % T3_wavepacket_full

\appendix
\relinput{.}{appendix/formule_explicite_archimedienne_FR} % explicit_arch_full

% --- section Passerelles interdisciplinaires
\relinput{.}{docs/explication-grand-format/latex/04_bridges_boxes_FR}

% --- section FAQ (par spécialité)
\relinput{.}{docs/explication-grand-format/latex/05_faq_personas_FR}

% --- section* Remerciements
\relinput{.}{docs/explication-grand-format/latex/remerciements_FR}

\bibliographystyle{alpha}
\relinput{.}{appendix/about_author_FR}
\bibliography{refs}
\end{document}
